\subsection{Ordinary ghosts for every fixed finite puncture set}

\begin{statement}{A.47}{Proved}{Simultaneous radical splitting with prescribed congruences}
Let $4\mid m$, let $d_1,\ldots,d_s$ be positive rational
numbers, and suppose there are distinct primes $P_i$ with
$v_{P_i}(d_i)\in\{1,-1\}$ and $v_{P_i}(d_j)=0$ for
$i\ne j$. Put $F=\Q(\mu_{2m})$ and suppose
$\sqrt{d_i}\in F$ for every $i$. Then
\[
 L=F(d_1^{1/m},\ldots,d_s^{1/m})
 \quad\text{satisfies}\quad
 L^{\mathrm{ab}}=F.
 \tag{A.47.1}
\]
For every multiple $M$ of $2m$ and every unit class
$u\pmod M$ with $u\equiv1\pmod {2m}$, infinitely many
primes split completely in $L$ and are congruent to $u$
modulo $M$.
\end{statement}
\begin{proof}
Choose the positive real roots $\alpha_i=d_i^{1/m}$ and
put $h=m/2$. We first determine all relations among the
classes of the $d_i$ in $F^\times/(F^\times)^m$.
If a positive rational $d$ has an $m$th root in the
abelian extension $F/\Q$, its positive real root
$\alpha$ also lies in $F$, since $\mu_m\subset F$.
The field $\Q(\alpha)$ is a subfield of $\R$ and is
Galois over $\Q$. All conjugates of $\alpha$ are
therefore real roots of $T^m-d$, and hence lie in
$\{\alpha,-\alpha\}$. Thus $\alpha^2\in\Q$, and
$d=(\alpha^2)^h$ is an $h$th power in $\Q$.

Apply this to $d=\prod_i d_i^{e_i}$. A relation
$d\in(F^\times)^m$ implies that each
$v_{P_i}(d)=\pm e_i$ is divisible by $h$.
Conversely $d_i^h=(\sqrt{d_i})^m$. The subgroup of
$F^\times/(F^\times)^m$ generated by these classes
is consequently $(\Z/h)^s$. The perfect Kummer
pairing
\[
 \operatorname{Gal}(L/F)\times
 \langle[d_1],\ldots,[d_s]\rangle
 \longrightarrow\mu_m,\qquad
 (\sigma,[d])\longmapsto
             \frac{\sigma(d^{1/m})}{d^{1/m}},
 \tag{A.47.2}
\]
gives $N:=\operatorname{Gal}(L/F)\simeq(\Z/h)^s$.
This is the explicit pairing of
\cite[Theorem 5.30 and Remark 5.32]{AgentMilneFT};
it follows from the power exact sequence and Hilbert 90.
We use coordinates in which the $i$th basis element
sends $\alpha_i$ to $\zeta_m^2\alpha_i$ and fixes
the other roots.

The extension $L/\Q$ is Galois because it contains all
roots of each rational polynomial $T^m-d_i$. Let
$G=\operatorname{Gal}(L/\Q)$, and let $\iota$ be
complex conjugation. It fixes the positive $\alpha_i$.
Conjugation by $\iota$ acts as $-1$ on $N$, so
$2N\subset[G,G]$. Also $[G,G]\subset N$, since
$F/\Q$ is abelian.

For $\sigma\in G$, write
$\sigma(\alpha_i)=\zeta_m^{k_i}\alpha_i$.
Its action on $\sqrt{d_i}=\alpha_i^h$ is
$(-1)^{k_i}$. Direct composition gives
\[
 [\iota,\sigma](\alpha_i)
                 =\zeta_m^{-2k_i}\alpha_i.
 \tag{A.47.3}
\]
In the coordinates of $N/2N$, this commutator is
$(-k_i\bmod2)_i$. The squareclasses of the $d_i$
in $\Q^\times/(\Q^\times)^2$ are independent by
the same prime valuations. The quadratic subfield
$\Q(\sqrt{d_1},\ldots,\sqrt{d_s})\subset F$ therefore
has Galois group $(\Z/2)^s$, and restriction from
$G$ realizes every vector of its quadratic
characters. Equation (A.47.3) shows that
$[G,G]$ surjects onto $N/2N$. Together with
$2N\subset[G,G]$, this gives $[G,G]=N$ and proves
(A.47.1).

For $2m\mid M$, the intersection
$L\cap\Q(\mu_M)$ equals $F$: it contains $F$
and is an abelian subextension of $L/\Q$.
The restriction map to the two fields identifies
the Galois group of their compositum with their
fibre product over $\operatorname{Gal}(F/\Q)$.
The pair consisting of the identity on $L$ and
$\zeta_M\mapsto\zeta_M^u$ belongs to this fibre
product. Chebotarev
\cite[Theorem 8.31]{AgentMilneANT} supplies the
claimed primes, outside any prescribed finite set.
\end{proof}

\begin{statement}{A.48}{Proved}{No fixed finite puncture set forces rationality}
For every finite set $S\subset\Q$ with at least two
elements there is a tuple
\[
 (x_v)_v\in
 \prod_v\bigl(\Q_v\setminus S\bigr)
\]
and classes $U_c\in\widehat{\Q^\times}$, one for
each $c\in S$, such that
\[
       (U_c)_p=\kappa(x_p-c)
       \quad\text{at every finite prime }p,
 \tag{A.48.1}
\]
with the corresponding equality of completed Kummer
classes at the real place. Every $U_c$ has infinitely
many nonzero ordinary integer valuation coordinates.
The tuple is not one rational number. In particular,
there is no finite $S\supset\{0,1\}$ for which the
finite form of (G) in A12 has only rational solutions,
even when every finite local coordinate is ordinary
and differs from every puncture.
\end{statement}
\begin{proof}
We give the construction for distinct integer
punctures $c_1,\ldots,c_k$, with $k\ge2$, and then
give the exact change of coordinates for rational
punctures.

\emph{Initial coefficients and fixed quadratic fields.}
Choose distinct primes $P_i>2$ larger than every
$|c_a-c_b|$ with $a\ne b$. By the Chinese remainder
theorem choose an integer $R>\max_i c_i$ such that
\[
                R\equiv c_i+P_i\pmod {P_i^2}
                  \quad(1\le i\le k).
 \tag{A.48.2}
\]
Put $a_{i,0}=R-c_i>0$. Then
$v_{P_i}(a_{i,0})=1$ and
$v_{P_i}(a_{t,0})=0$ for $t\ne i$; also no $P_i$
divides a nonzero puncture difference.
Let
\[
 D_0=\prod_i a_{i,0}
             \prod_{i<t}|c_i-c_t|,\qquad
 N_0=8\,\operatorname{rad}(D_0),\qquad
 n_j=N_0(j+1)!.
 \tag{A.48.3}
\]
The field $\Q(\mu_{N_0})$ contains every
$\sqrt{|(c_i-c_t)/a_{t,0}|}$. Here is an explicit
verification. For an odd prime $P$, the quadratic
Gauss sum belongs to $\Q(\mu_P)$ and its square
is $(-1/P)P$: its squared absolute value is $P$,
as computed in A.44, and complex conjugation
multiplies it by $(-1/P)$. Since $i\in\Q(\mu_{N_0})$,
this puts $\sqrt P$ in that field for each odd
$P\mid D_0$. Also
$\sqrt2=\zeta_8+\zeta_8^{-1}$ lies there.
Factoring the numerator and denominator into prime
powers proves the assertion.

\emph{A uniform large-prime bound for the Kummer curves.}
For an exponent $m$ prime to $p$, distinct residues
$c_i\pmod p$, and units $a_i\pmod p$, consider the
affine curve
\[
          T-c_i=a_iX_i^m\quad(1\le i\le k),
          \qquad \prod_i X_i\ne0.
 \tag{A.48.4}
\]
It is a smooth geometrically connected cover of
$\mathbb A^1_{\mathbf F_p}\setminus\{c_i\}$ of
degree $m^k$. Indeed, over an algebraic closure,
the classes of $T-c_i$ in
$\overline{\mathbf F}_p(T)^\times/
       (\overline{\mathbf F}_p(T)^\times)^m$
are independent, as their valuations at $c_i$
are the standard basis vectors. The Kummer
pairing (A.47.2) gives degree $m^k$, and the
coordinate derivatives in $X_i$ are units on
the displayed affine curve.

Its smooth projective completion has ramification
index $m$ above each $c_i$ and above infinity,
and no other ramification. At $c_i$, precisely
the corresponding Kummer function has valuation
one; all other functions are units with $m$th
roots in the strict local ring. At infinity,
with parameter $z=1/T$, each function is a
constant times $z^{-1}(1-c_i z)$, so a single
$m$th root of $z$ accounts for the ramification.
The tame Riemann--Hurwitz formula consequently gives
\[
 2g-2=(k-1)m^k-(k+1)m^{k-1}.
 \tag{A.48.5}
\]
The removed fibres contain exactly
$(k+1)m^{k-1}$ geometric points.
The Hasse--Weil bound therefore makes the number
of interior $\mathbf F_p$-points at least
\[
 p+1-2g\sqrt p-(k+1)m^{k-1}.
 \tag{A.48.6}
\]
For the precise cohomological statement used here,
the trace formula is \cite[(1.5.1)]{AgentDeligneWeilI},
the absolute values of the eigenvalues in degree
one are \cite[Theorem (1.6)]{AgentDeligneWeilI},
and that degree has dimension $2g$ by
\cite[Proposition 14.2]{AgentMilneLEC}.

Set
\[
             B_k(m)=\bigl((k+1)m^k+3\bigr)^2.
 \tag{A.48.7}
\]
For $T_0=(k+1)m^k+3$, (A.48.5) gives
$T_0-2g=2m^k+(k+1)m^{k-1}+1$.
Thus the quadratic
$T^2-2gT+1-(k+1)m^{k-1}$ is positive at $T_0$
and increasing for $T\ge T_0$.
Consequently (A.48.6) is positive whenever
$p>B_k(m)$. An interior residue point lifts
to a point with all $X_i\in\Z_p^\times$:
keep a lift of $T$ fixed and apply simple-root
Hensel lifting to each of the $k$ equations.

\emph{The simultaneous induction.}
We construct positive integer coefficients $a_{i,j}$
with private valuations $v_{P_i}(a_{i,j})=1$ and
$v_{P_i}(a_{t,j})=0$ for $t\ne i$. Each coefficient
is its initial value times eighth powers of new
primes. At every finite prime $p$ we require a
point
\[
 T_p-c_i=a_{i,j}X_{i,p}^{n_j},
       \qquad X_{i,p}\in\Z_p^\times
       \quad(1\le i\le k).
 \tag{A.48.8}
\]
Put $V_j(p)=\max_i v_p(a_{i,j})$ and require also
\[
 \begin{split}
 v_p\bigl((a_{i,j}+c_i)-(a_{1,j}+c_1)\bigr)
       &\ge H_j(p),\\
 H_j(p)&=V_j(p)+2v_p(2n_{j+1})+j+3,
       \qquad p\le B_k(n_j).
 \end{split}
 \tag{A.48.9}
\]
The initial coefficients satisfy both conditions
with $T_p=R$ and all $X_{i,p}=1$.

At stage $j$, write $n=n_j$, $m=n_{j+1}$,
and $a_i=a_{i,j}$. Let $\mathcal T$ consist of
the primes at most $B_k(m)$, the prime divisors
of the $a_i$, and the prime divisors of the
nonzero differences $c_i-c_t$. For $p\in\mathcal T$
put
\[
 L_p=V_j(p)+2v_p(2n_{j+2})+j+4,\qquad
 M=\prod_{p\in\mathcal T}p^{L_p}.
 \tag{A.48.10}
\]
This set contains all prime divisors of $2m$,
and $2m\mid M$.

Choose a local point (A.48.8). At $p\mid2m$
choose it with
\[
                 X_{i,p}\equiv1\pmod {p^{v_p(2m)}}
                       \quad\text{for every }i.
 \tag{A.48.11}
\]
To justify this choice, every such $p$ divides
$2N_0$ or is at most $j+2$, and hence
$p\le B_k(n)$. Set $T_p=a_1+c_1$ and
$X_{1,p}=1$. For $i\ne1$ the quotient
\[
 \frac{T_p-c_i}{a_i}
    =1+\frac{(a_1+c_1)-(a_i+c_i)}{a_i}
\]
differs from one by an element of valuation at
least $2v_p(2m)+j+3$. Applying (A.44.7) to its
$n$th root gives a unit root whose difference
from one has valuation at least
$2v_p(2m)+j+3-v_p(n)\ge v_p(2m)$.
This proves (A.48.11), including at two and at
primes dividing several coefficients.

We now choose fresh primes $q_1,\ldots,q_k$
sequentially. Suppose $q_1,\ldots,q_{i-1}$
have been chosen. Set
\[
 b_t=\begin{cases}
          a_tq_t^n,&t<i,\\
          a_t,&t\ge i,
       \end{cases}
 \qquad
 D_{it}=\frac{c_i-c_t}{b_t}\quad(t\ne i).
 \tag{A.48.12}
\]
Consider the Galois field
\[
 F=\Q(\mu_{2m}),\qquad
 L_i=F\bigl(|D_{it}|^{1/m}:t\ne i\bigr).
 \tag{A.48.13}
\]
It contains $m$th roots of the signed $D_{it}$
as well, since an $m$th root of $-1$ lies in $F$.
For each $t\ne i$, the private prime $P_t$
has valuation $-1$ on $|D_{it}|$ and valuation
zero on all other radicands of (A.48.13).
The numerator differences are $P_t$-units by
(A.48.2); prior updates use different fresh
primes and do not change these private valuations.
Moreover the squareclass of $|D_{it}|$ is the
fixed squareclass of $|(c_i-c_t)/a_{t,0}|$,
because every update exponent is even.
The square roots of these numbers already lie
in $\Q(\mu_{N_0})\subset F$. Statement A.47
therefore gives
\[
                        L_i^{\mathrm{ab}}=F.
 \tag{A.48.14}
\]

At each previously selected $q_h$, with $h<i$,
there is a nonzero residue $\rho_{hi}$ satisfying
\[
              \rho_{hi}^{\,n}
                    \equiv(c_h-c_i)/a_i\pmod {q_h}.
 \tag{A.48.15}
\]
Indeed, the choice of $q_h$ required complete
splitting in the field (A.48.13) for index $h$;
at that time coordinate $i$ had not yet been
updated. That field contains an $m$th root of
$(c_h-c_i)/a_i$, and $n\mid m$.
Now impose on $q_i$ the finite congruences
\[
 \begin{aligned}
 q_i&\equiv X_{i,p}\pmod {p^{L_p}}
                         &&(p\in\mathcal T),\\
 q_i&\equiv\rho_{hi}\pmod {q_h}
                         &&(h<i).
 \end{aligned}
 \tag{A.48.16}
\]
They specify a unit residue class modulo
$M\prod_{h<i}q_h$. By (A.48.11), its reduction
modulo $2m$ is one. Each earlier $q_h$ is prime
to $2m$ and to $M$, so its prescribed residue
in the second line imposes no additional
condition modulo $2m$. Applying A.47 with
(A.48.14), choose $q_i$ in this residue class
which splits completely in $L_i$ and is larger
than all primes excluded or previously chosen
and than $n_{j+2}$. There are infinitely many
choices. This completes all $k$ selections.

Define
\[
                 a_{i,j+1}=a_iq_i^n,\qquad
                 q_{i,j}=q_i.
 \tag{A.48.17}
\]
All the new primes avoid the supports of the
old coefficients and the puncture differences.
The private valuations remain as required,
and the coefficients retain their initial
squareclasses and eighth-power forms.

We verify the new local condition at every prime.
For $p\in\mathcal T$, (A.48.16) gives
$a_iq_i^n+c_i\equiv T_p\pmod {p^{L_p}}$
for every $i$. Thus all their pairwise
differences have valuation at least $L_p$.
Also $v_p(a_iq_i^n)=v_p(a_i)$, since $q_i\ne p$.
Set $T'_p=a_{1,j+1}+c_1$. Each quotient
$(T'_p-c_i)/a_{i,j+1}$ differs from one by
an element of valuation at least
$L_p-V_j(p)>2v_p(m)$. Hensel's estimate gives
an $m$th root in $\Z_p^\times$, proving the
new local condition. Since $V_{j+1}(p)=V_j(p)$
on $\mathcal T$, (A.48.10) is exactly the next
precision bound in (A.48.9) for every
$p\le B_k(m)$.

At a new prime $q_h$, set $X_h=1$ and
$T=c_h+a_hq_h^n$. For $i<h$, the quotient
$(T-c_i)/(a_iq_i^n)$ reduces modulo $q_h$
to $D_{hi}$ from the field used when $q_h$
was chosen, and this residue is an $m$th
power by complete splitting. For $i>h$,
(A.48.15)--(A.48.16) give
$a_iq_i^n\equiv c_h-c_i\pmod {q_h}$,
so the quotient reduces to one. The numerators
and denominators in these quotients are units,
because $q_h$ divides no puncture difference.
As $q_h\nmid m$, all these $m$th roots lift
to $\Z_{q_h}^\times$. This verifies the
simultaneous local point at every new prime.

At every remaining prime the coefficients and
puncture differences are units, the exponent
$m$ is invertible, and the prime exceeds
$B_k(m)$. Equations (A.48.4)--(A.48.7)
give an interior residue point and hence a
unit-coordinate local point. All finite
primes are now covered, completing the induction.

\emph{The limits and the rational-puncture comparison.}
For each $i$, the product
\[
                  U_i=a_{i,0}
                         \prod_{j\ge0}q_{i,j}^{\,n_j}
                     \quad\text{in }\widehat{\Q^\times}
 \tag{A.48.18}
\]
is defined because the $n_j$ are eventually
divisible by every positive integer. All selected
primes are distinct throughout the construction.
Their valuations in coordinate $i$ are $n_j$,
and their valuations in the other coordinates
are zero. The valuation at every fixed prime
is eventually constant in every coefficient
sequence. As in A.44, factorial divisibility
annihilates each finite local-unit quotient,
so $a_{i,j}$ converges in $\Q_p^\times$ to
an ordinary nonzero $y_{i,p}$ with
$(U_i)_p=\kappa(y_{i,p})$.
Equation (A.48.9) shows that
$y_{i,p}+c_i$ is independent of $i$.
Call its common value $x_p$. Then
$x_p-c_i=y_{i,p}\ne0$ for every puncture,
and (A.48.1) holds. At the real place all
$U_i$ have positive sign, represented by
any $x_\infty>\max_i c_i$.

Every $U_i$ has an infinite set of positive
valuation coordinates at its own selected
primes. The tuple cannot equal a rational
number $r$: that would force the nonzero
rational $r-c_i$ to have these infinitely
many nonzero valuations.

Finally, for distinct rational punctures
$c_i$, choose a positive integer $D$ with
$C_i=Dc_i\in\Z$. The $\Q$-isomorphism
\[
 f_D:\mathbb P^1_\Q\setminus(S\cup\{\infty\})
       \longrightarrow
       \mathbb P^1_\Q\setminus(\{C_i\}\cup\{\infty\}),
       \qquad t\longmapsto Dt
 \tag{A.48.19}
\]
has $f_D^*(T-C_i)=D(t-c_i)$. Apply the integer
construction to obtain $x'_v$ and $V_i$, and put
$x_v=x'_v/D$ and $U_{c_i}=V_i\kappa(D)^{-1}$.
Multiplication by $\kappa(D)^{-1}$ is a continuous
translation of the completed multiplicative group.
Localizing the displayed identity proves
$ (U_{c_i})_p=\kappa(x_p-c_i)$.
Division by $D$ changes only finitely many
valuation coordinates, so each $U_{c_i}$ still
has infinite support. The positive real sign
is preserved. This gives all assertions over
the original rational puncture set.
\end{proof}

The construction in A.48 is made for each fixed
finite $S$. It does not supply one tuple together
with compatible classes for every $c\in\Q$;
the prime choices and the initial quadratic
fields depend on $S$. The countably punctured
ordinary finite-place form of (G) therefore
requires an additional construction or obstruction.

The accompanying exact certificate
\texttt{code/finite\_puncture\_certificate.py}
checks the commutator identity, private-prime
valuations after sequential updates, the genus
and boundary formulas, the large-prime inequality,
the simultaneous roots at a new coefficient prime,
and the $V_j(p)$-corrected precision inequalities.
Its saved output records all finite checks.
It does not execute a full prime-search stage;
termination of those searches is established
by A.47 and Chebotarev in the proof.
